% Generated by scripts/build_reports.py; edit reviewed Markdown inputs.
% Compile with XeLaTeX.
% Options for packages loaded elsewhere
\PassOptionsToPackage{unicode}{hyperref}
\PassOptionsToPackage{hyphens}{url}
\PassOptionsToPackage{dvipsnames,svgnames,x11names}{xcolor}
\documentclass[
  11pt,
]{article}
\usepackage{xcolor}
\usepackage[margin=1in]{geometry}
\usepackage{amsmath,amssymb}
\setcounter{secnumdepth}{-\maxdimen} % remove section numbering
\usepackage{iftex}
\ifPDFTeX
  \usepackage[T1]{fontenc}
  \usepackage[utf8]{inputenc}
  \usepackage{textcomp} % provide euro and other symbols
\else % if luatex or xetex
  \usepackage{unicode-math} % this also loads fontspec
  \defaultfontfeatures{Scale=MatchLowercase}
  \defaultfontfeatures[\rmfamily]{Ligatures=TeX,Scale=1}
\fi
\usepackage{lmodern}
\ifPDFTeX\else
  % xetex/luatex font selection
\fi
% Use upquote if available, for straight quotes in verbatim environments
\IfFileExists{upquote.sty}{\usepackage{upquote}}{}
\IfFileExists{microtype.sty}{% use microtype if available
  \usepackage[]{microtype}
  \UseMicrotypeSet[protrusion]{basicmath} % disable protrusion for tt fonts
}{}
\makeatletter
\@ifundefined{KOMAClassName}{% if non-KOMA class
  \IfFileExists{parskip.sty}{%
    \usepackage{parskip}
  }{% else
    \setlength{\parindent}{0pt}
    \setlength{\parskip}{6pt plus 2pt minus 1pt}}
}{% if KOMA class
  \KOMAoptions{parskip=half}}
\makeatother
\usepackage{longtable,booktabs,array}
\usepackage{caption}
\captionsetup[table]{skip=6pt}
\newcounter{none} % for unnumbered tables
\usepackage{calc} % for calculating minipage widths
% Correct order of tables after \paragraph or \subparagraph
\usepackage{etoolbox}
\makeatletter
\patchcmd\longtable{\par}{\if@noskipsec\mbox{}\fi\par}{}{}
\makeatother
% Allow footnotes in longtable head/foot
\IfFileExists{footnotehyper.sty}{\usepackage{footnotehyper}}{\usepackage{footnote}}
\makesavenoteenv{longtable}
\setlength{\emergencystretch}{3em} % prevent overfull lines
\providecommand{\tightlist}{%
  \setlength{\itemsep}{0pt}\setlength{\parskip}{0pt}}
\usepackage{amsmath,amssymb}
\usepackage{xurl}
\usepackage{fancyhdr}
\pagestyle{fancy}
\fancyhf{}
\fancyhead[L]{\small Pure Thought | Revised research specifications}
\fancyfoot[L]{\small 2026-09-17}
\fancyfoot[R]{\small\thepage}
\setlength{\headheight}{15pt}
\setlength{\emergencystretch}{3em}
\allowdisplaybreaks
\usepackage{bookmark}
\IfFileExists{xurl.sty}{\usepackage{xurl}}{} % add URL line breaks if available
\urlstyle{same}
% fallback for those not using the hyperref driver hyperxmp:
\makeatletter
\@ifundefined{xmpquote}{\newcommand{\xmpquote}[1]{#1}}{}
\makeatother
\hypersetup{
  pdftitle={18. Certified optimal-transport bounds for strictly correlated electrons},
  pdfauthor={Pure Thought research portfolio},
  colorlinks=true,
  linkcolor={blue},
  filecolor={Maroon},
  citecolor={Blue},
  urlcolor={blue},
  pdfcreator={LaTeX via pandoc}}

\title{18. Certified optimal-transport bounds for strictly correlated
electrons}
\author{Pure Thought research portfolio}
\date{Revised 17 September 2026}

\begin{document}
\maketitle
\thispagestyle{empty}
\clearpage

\section{18. Certified optimal-transport bounds for strictly correlated
electrons}

\textbf{Revised specification --- 2026-09-11.}
\href{https://github.com/rick-stevens-ai/Pure-Thought/blob/main/reviews/18-Optimal-Transport-Chemistry.md}{Detailed
critique} \ensuremath{\cdot}
\href{https://github.com/rick-stevens-ai/Pure-Thought/blob/main/archive/original-PRDs/18-Optimal-Transport-Chemistry.md}{Original
description} \ensuremath{\cdot}
\href{https://github.com/rick-stevens-ai/Pure-Thought/blob/main/reviews/EVIDENCE-STANDARD.md}{Evidence
standard}

\subsection{Plain-language guide}

\subsubsection{The problem in everyday terms}

Electrons repel one another, but a prescribed electron density also
constrains where they spend their time. Imagine arranging their joint
positions so they avoid one another as effectively as possible while
preserving that average density. This problem makes that arrangement
question precise and seeks reliable upper and lower bounds on its
repulsion energy. It represents a limiting strong-correlation model, not
the full energy of ordinary electrons.

\subsubsection{Key terms}

\begin{itemize}
\item
  \textbf{Electron density \ensuremath{\rho}} --- The average spatial
  distribution of electrons, integrating to their total number N.
\item
  \textbf{Coulomb cost} --- The repulsion contribution between charged
  particles, proportional here to the inverse of their separation.
\item
  \textbf{Optimal transport} --- Optimization over ways of distributing
  or coupling quantities while satisfying prescribed distribution
  constraints.
\item
  \textbf{Marginal distribution} --- The distribution of one variable
  after the other variables in a joint distribution are ignored.
\item
  \textbf{Coupling \ensuremath{\Gamma}} --- A joint distribution of all
  electron positions consistent with the prescribed individual
  marginals.
\item
  \textbf{Multi-marginal transport} --- Transport optimization with more
  than two prescribed marginals.
\item
  \textbf{Strictly correlated electrons (SCE)} --- A limiting model in
  which electron positions are organized to minimize interaction energy
  at fixed density.
\item
  \textbf{Dual potential} --- An auxiliary function providing a
  certified lower bound when it satisfies the required inequality
  everywhere.
\item
  \textbf{Discretization error} --- The difference introduced by
  replacing a continuous spatial problem with a finite representation.
\end{itemize}

\subsubsection{Why this matters}

Strongly correlated electrons are difficult to approximate because their
motions are strongly linked. This limiting problem gives a
mathematically structured reference for studying that linkage and for
developing density-functional ideas. Verified bounds help distinguish a
true property of the continuum model from an artifact of a spatial grid
or optimization tolerance.

\subsubsection{What progress would mean}

A useful result would enclose the SCE interaction energy for a declared
density. It could inform further electronic-structure theory, while
leaving kinetic energy, finite-correlation effects and actual molecular
reaction paths to additional modeling.

\subsection{Core question}

For a supplied normalized electron density, can a discretized Coulomb
multi-marginal transport problem yield verified bounds, with continuum
error separated from optimization error?

\subsection{Scope and assumptions}

Take an electron density \ensuremath{\rho}\ensuremath{\ge}0 with
\ensuremath{\int}\ensuremath{\rho}=N and marginal
\ensuremath{\mu}=\ensuremath{\rho}/N. Begin with N=2 and a compactly
supported mathematical density or a finite rational discrete model.
Specify treatment of the diagonal Coulomb singularity. A
continuous-density result additionally needs domain-tail and
discretization estimates.

The main target is the strictly correlated interaction functional.
Wasserstein comparison of densities is a separate optional metric task;
it is not a minimum-energy chemical reaction path.

\subsection{Mathematical target}

Define

\[V_{SCE}[\rho ]=\inf_{\Gamma \in\Pi(\mu ,\ldots,\mu )}\int\sum_{i<j}|r_i-r_j|^{-1}\,d\Gamma .\]

\ensuremath{\Gamma} is a probability measure and may be symmetrized; it
is not an antisymmetric probability distribution. The symmetric dual
uses u with \ensuremath{\Sigma}\_i
u(r\_i)\ensuremath{\le}\ensuremath{\Sigma}\_\{i<j\}|r\_i\ensuremath{-}r\_j|\ensuremath{{}^{-1}}
and objective \ensuremath{\int}u\ensuremath{\rho}. A feasible coupling
gives an upper bound and a dual-feasible potential gives a lower bound
for the same problem. Define
W\_\ensuremath{\infty}=V\_SCE\ensuremath{-}U\_H if subtracting Hartree
energy; do not call V\_SCE exact exchange.

\subsection{Required outputs and proof obligations}

Return marginal conventions, cost and singularity policy, primal plan,
dual potentials, rigorously checked constraints, and
L\ensuremath{\le}V\ensuremath{\le}U. State whether V is a finite
discrete optimum or the continuum functional. Sinkhorn candidates must
be converted to feasible bounds for the unregularized problem, with
entropic bias controlled.

\subsection{Validation and rejection controls}

Verify a small rational-cost transport LP exactly and a one-dimensional
quadratic-cost quantile example as a separate solver test. For Coulomb
densities, use a case with a known co-motion construction under its
actual hypotheses. Reject comparing normalized densities with unequal
total mass unless normalization or an unbalanced-OT model is explicitly
chosen.

\subsection{Milestones and research extension}

Baseline: certified two-marginal discrete bounds. Strong: one continuum
enclosure with singularity/tail control. Research extension:
N\textgreater2 with symmetry reduction and validated bounds, or a
demonstrably improved relaxation. Do not promise deterministic Monge
maps for arbitrary multimarginal inputs.

Milestones are acceptance gates, not calendar promises. Before a run,
freeze the concrete instance/domain, tolerance, compute budget, and
permitted dependencies. A bounded search that fails to find a result
must preserve an unresolved outcome.

\subsection{Starting instruction}

Work on the baseline above. First make every input and convention
explicit, then validate the reference case. Build the checker alongside
the calculation. Label numerical evidence, finite-model certificates,
and continuum theorems separately; return the strongest justified result
and identify the exact unmet proof obligation.

\subsection{Primary references and their role}

\begin{itemize}
\tightlist
\item
  \href{https://arxiv.org/abs/1104.0603}{Cotar, Friesecke and
  Klüppelberg, Density functional theory and optimal transportation with
  Coulomb cost} --- Strong-correlation/semiclassical connection; theorem
  hypotheses must be checked for each density.
\end{itemize}

References support definitions or limitations; the scope choices and
proposed milestones are editorial recommendations, not claims that these
results have already been obtained.

\subsection{Displayed formulation and claim boundary}

\href{https://github.com/rick-stevens-ai/Pure-Thought/blob/main/PRDs/18-Optimal-Transport-Chemistry.md}{Canonical
specification}.

\[\int\rho=N,\qquad \mu=\rho/N,\qquad
c(r_1,\ldots,r_N)=\sum_{i<j}|r_i-r_j|^{-1},\]

\[V_{\mathrm{SCE}}[\rho]=\inf_{\Gamma\in\Pi(\mu,\ldots,\mu)}\int c\,d\Gamma,
\qquad \sum_i u(r_i)\le c(r_1,\ldots,r_N),\]

\[\underbrace{\int u\rho}_{L}\ \le\ V_{\mathrm{SCE}}[\rho]\ \le\
\underbrace{\int c\,d\Gamma_{\mathrm{feasible}}}_{U}.\]

The coupling is a probability measure; it may be symmetrized and is not
antisymmetric. This is a Coulomb interaction functional, not exact
exchange. For a discrete problem, the same inequalities apply only to
the discrete optimum until continuum error is bounded.

\subsection{Detailed review and rationale}

\href{https://github.com/rick-stevens-ai/Pure-Thought/blob/main/PRDs/18-Optimal-Transport-Chemistry.md}{Revised
problem} \ensuremath{\cdot}
\href{https://github.com/rick-stevens-ai/Pure-Thought/blob/main/archive/original-PRDs/18-Optimal-Transport-Chemistry.md}{Preserved
original} \ensuremath{\cdot}
\href{https://github.com/rick-stevens-ai/Pure-Thought/blob/f854d4e12b55c1f5b5c2a964b75f2d23496d981e/PRDs/18-Optimal-Transport-Chemistry.md}{Pinned
upstream source}

Line references below refer to the preserved original at commit
\texttt{f854d4e12b55}, not the rewritten file. ``Blocking'' means the
original issue can invalidate the scientific target or claimed
certificate; ``major'' means it materially changes scope, inference, or
acceptance criteria. These are critiques of specifications, not results
of running the embedded research-code sketches.

\subsubsection{Detailed findings}

\begin{enumerate}
\def\labelenumi{\arabic{enumi}.}
\item
  \textbf{Blocking --- wrong functional (16, 34, 98--108).} Coulomb OT
  describes a strong-correlation limit, not the exact exchange
  functional asserted here. The negative half-cost formula and
  negative-cost dual are inappropriate for the positive repulsion
  minimization.
\item
  \textbf{Major --- normalization and antisymmetry.} The electron
  density integrates to N, while each probability marginal integrates to
  one. Fermionic wavefunction antisymmetry does not make a probability
  measure antisymmetric.
\item
  \textbf{Major --- reaction pathway inference (18, 36).} A Wasserstein
  geodesic minimizes transport cost, not a molecular potential-energy
  barrier. The original chemical interpretation lacks a
  dynamical/energetic bridge.
\item
  \textbf{Major --- uniqueness and certification (106, 868--905).}
  Quadratic-cost Brenier results do not transfer automatically to
  Coulomb multimarginal problems. Marginal residuals and a regularized
  solver output are not continuum certificates. The revised formulation
  distinguishes these errors and removes the inappropriate
  helium-exchange target.
\end{enumerate}

\subsubsection{What the revision preserves and changes}

The central motivation is retained, but the executable question is now:
\textbf{For a supplied normalized electron density, can a discretized
Coulomb multi-marginal transport problem yield verified bounds, with
continuum error separated from optimization error?} The revised scope
chooses a concrete starting model; those choices are proposed
research-design decisions, not facts inferred from the original. The
baseline can succeed without resolving the full open research program.

The replacement also removes the original implementation sketches from
the active specification. They remain in the archive for comparison, but
have not been made executable or certified. New implementations should
derive their equations from the revised definitions rather than copy
unchecked prototypes.

\subsubsection{Acceptance criteria to use instead}

Baseline: certified two-marginal discrete bounds. Strong: one continuum
enclosure with singularity/tail control. Research extension:
N\textgreater2 with symmetry reduction and validated bounds, or a
demonstrably improved relaxation. Do not promise deterministic Monge
maps for arbitrary multimarginal inputs.

\subsubsection{Sources supporting the corrections}

\begin{itemize}
\tightlist
\item
  \href{https://arxiv.org/abs/1104.0603}{Cotar, Friesecke and
  Klüppelberg, Density functional theory and optimal transportation with
  Coulomb cost} --- Strong-correlation/semiclassical connection; theorem
  hypotheses must be checked for each density.
\end{itemize}

\section{Evidence and execution standard}

This standard applies to the revised specifications. It is a proposed
contract for future research runs; this editorial revision does not
supply the scientific certificates themselves.

\subsection{What ``pure thought'' means here}

Use explicitly stated mathematical models and deductions without fitting
conclusions to experimental outcomes. Synthetic examples, exact model
inputs, analytic references, numerical search, and established software
are allowed. Physical constants, electronic-structure approximations,
constitutive laws, and empirical parameter sets must be identified when
used. A model-derived result is conditional on that model.

The original ban on legacy software conflicts with its proposed reliance
on NumPy, SciPy, CVXPY, nauty, and other libraries. Reimplementing every
optimizer and algebra system is not necessary for scientific
independence. Prefer a small independently checkable certificate and a
documented trusted base. A fresh implementation can be a project
objective, but must not be confused with a mathematical proof or novelty
claim.

\subsection{Four evidence levels}

{\def\LTcaptype{none} % do not increment counter
\begin{longtable}[]{@{}
  >{\raggedright\arraybackslash}p{(\linewidth - 4\tabcolsep) * \real{0.3333}}
  >{\raggedright\arraybackslash}p{(\linewidth - 4\tabcolsep) * \real{0.3333}}
  >{\raggedright\arraybackslash}p{(\linewidth - 4\tabcolsep) * \real{0.3333}}@{}}
\toprule\noalign{}
\begin{minipage}[b]{\linewidth}\raggedright
Label
\end{minipage} & \begin{minipage}[b]{\linewidth}\raggedright
What it establishes
\end{minipage} & \begin{minipage}[b]{\linewidth}\raggedright
What it does not establish
\end{minipage} \\
\midrule\noalign{}
\endhead
\bottomrule\noalign{}
\endlastfoot
Numerical evidence & A reproducible computation at stated precision,
sample size, and discretization & Exact truth, continuum coverage, or a
universal theorem \\
Finite-model certificate & An exact/validated statement about the
specified finite model or finite relaxation & An omitted tail,
thermodynamic limit, or physical realization \\
Conditional mathematical theorem & A proof over the stated domain with
every hypothesis and error controlled & Truth outside those assumptions
or a demonstrated physical system \\
Formal proof & A checked theorem in a named proof assistant, with
imported axioms/dependencies listed & Correct physical modeling merely
by virtue of formalization \\
\end{longtable}
}

Formalization is an optional extra deliverable unless the specific
project makes it mandatory. High precision, JSON output, solver status
and small residuals do not by themselves confer any of these proof
levels.

\subsection{Claim directions}

For minimization over a larger feasible set, the optimum is a lower
bound on the original minimum. A physical feasible point supplies an
upper bound. For maximization the directions reverse. Track feasibility
of the primal/dual witnesses separately from their objective values.

Infeasibility of an inner ansatz excludes only that ansatz.
Infeasibility of an outer relaxation can exclude the original problem if
the map into the relaxation is proved and all approximations preserve
necessity. Passing necessary tests means ``not excluded''; existence
needs a constructive sufficient argument.

A smallest observed example is not a minimality theorem. A complete
catalog must define the universe, equivalence relation, boundaries, and
proof that pruning misses nothing. A finite-size numerical trend does
not prove an asymptotic limit.

\subsection{Minimum run manifest}

Before calculation, record the problem ID and revision hash, exact
inputs and conventions, quantified domain, equivalence relation, target
observable, accepted evidence level, tolerances, and resource limits.
Record dependencies and precision/rounding mode. Where randomization is
used, preserve seeds, trial counts and raw sufficient statistics.
Resource limits are planning decisions, not evidence that a mathematical
solution cannot exist.

An output bundle should include:

\begin{itemize}
\tightlist
\item
  The statement being claimed, including all assumptions and
  quantifiers.
\item
  Construction or proof object, plus an independently runnable verifier
  and its trust assumptions.
\item
  Numerical/error ledger separating roundoff, discretization,
  truncation, optimization, sampling and modeling error where relevant.
\item
  Positive reference cases and negative controls that the verifier must
  reject.
\item
  Coverage: included domain, omitted domain, unresolved regions and
  termination reason.
\end{itemize}

These are requirements for future implementations, not a promise of
polynomial-size certificates for every problem.

\subsection{Checking specific certificate types}

\textbf{LP/SDP/SOS:} verify equality constraints, objective direction,
and cone feasibility using exact arithmetic or outward-rounded
intervals. Rationalizing a numerical solution can break PSD constraints.
Correct residuals with justified norm bounds; do not silently project
and reuse the old objective. Distinguish a proof of a polynomial
approximation from a proof for the original function.

\textbf{Interval numerics:} state function enclosures, domains and
outward rounding. A small root box needs existence and, if claimed,
uniqueness. Cover the complement for a global enumeration. Include
singularities and tails rather than discarding them numerically.

\textbf{SAT/SMT:} a solver's UNSAT label is not a certificate. Preserve
a proof supported by a checker and the proof that the encoding, symmetry
reduction and blocked solutions match the mathematical universe. DRAT
applies to an appropriate Boolean CNF proof, not automatically to every
real-arithmetic or SDP claim.

\textbf{Statistics:} specify the estimand, dependence assumptions,
confidence procedure and bias. Zero observed failures does not prove
zero failure probability. MCMC error bars based on estimated
autocorrelation are not generally rigorous mixing guarantees.

\subsection{Milestones and novelty}

Baseline means a working, independently validated result on a concrete
reference problem. Strong means a broader or sharper result with its
claimed evidence level. Research extension means a possible novel result
whose feasibility and significance need review. A negative theorem can
be valuable; an unsuccessful search can also be useful if its scope is
documented, but the two are not interchangeable.

Do not promise ``publication-quality'' in a fixed number of months.
Before claiming novelty, conduct a focused current literature comparison
against the exact theorem/model/metric. The sources in this revision
check key definitions and limitations; they are not an exhaustive
novelty audit across all 30 fields.

\end{document}
